\chapter{Manual de usuario}

Este capítulo explica cómo utilizar la aplicación una vez desplegada. Dado que el sistema cuenta con
un único rol de usuario, el manual se organiza en función de las tareas que este puede realizar, en
lugar de por perfiles. El contenido se irá ampliando a medida que se completen los sucesivos
\textit{sprints}.

\section{Acceso y visión general}

Al abrir la aplicación en el navegador se muestra una pantalla de acceso de carácter presentacional.
Tras acceder, se presenta la red cerebral de ejemplo (atlas AAL90) ya renderizada en el lienzo 3D,
junto con los paneles de información y control superpuestos.

% NOTA: incluir captura de la pantalla de acceso y de la vista general.

\section{Navegación por la escena 3D}

El usuario puede explorar la red mediante los controles de cámara:

\begin{itemize}
	\item \textbf{Rotar (orbitar):} arrastrar con el botón izquierdo del ratón.
	\item \textbf{Zoom:} rueda del ratón o gesto de pellizco en dispositivos táctiles.
	\item \textbf{Desplazar (paneo):} arrastrar con el botón derecho o con dos dedos.
\end{itemize}

\section{Interpretación de la visualización}

Cada nodo representa una región cerebral situada en sus coordenadas anatómicas, y su color indica el
grupo o lóbulo al que pertenece, según la leyenda mostrada en la interfaz. Las líneas entre nodos
representan las conexiones de la red. El panel de estadísticas muestra el número de nodos,
conexiones y grupos de la red activa.

% NOTA: a medida que avancen los sprints, documentar aquí el filtrado por umbral (RF 5), la
% selección e inspección de nodos (RF 6, RF 6.1), la consulta de métricas (RF 8), la búsqueda
% (RF 10), y las funciones de investigador/administrador: registro (RF 13), carga de datasets
% propios (RF 14), gestión de datasets (RF 15) y exportación (RF 12, RF 16).
